% SOURCE_AUTHORITY: sga5_fr_workpass.tex, lines 12224--12236 (LF-normalized unit).
% SOURCE_SHA256: 791F4EFFC5E02832D5D77ED03518C8156D6F07E4C8238B03545DB93D883FBB28
\section*{Exposição X --- Fórmula de Euler--Poincaré em cohomologia étale}

\begin{center}
\textbf{FÓRMULA DE EULER--POINCARÉ EM COHOMOLOGIA ÉTALE}

por A. Grothendieck

(redigido por I. Bucur)
\end{center}

O objeto da presente exposição é o enunciado e a demonstração da fórmula (7.2), do tipo Euler--Poincaré. Esta fórmula generaliza uma fórmula obtida independentemente por Ogg e Chafarevitch, cujos resultados são expostos (do ponto de vista da cohomologia étale) na exposição de M. Raynaud [1].

A fórmula de Euler--Poincaré utiliza de modo essencial o formalismo das categorias derivadas. Como em loc. cit., a demonstração é efetuada por meio de uma variante (5.1) de uma conhecida fórmula de Weil.
